\scanpage{171}
\textcircled{3} L'opération de Cartier.

\uline{Théorème 1} Supposons que $B/A$ admette une $p$-base. Alors l'hom précédent $\varphi$ \deleted{est}\inserted{induit} un \uline{isomorphisme} $\Omega^*_{\widetilde B/\widetilde A}\otimes_{\widetilde B}\Lambda\longrightarrow H^{\unsure{\alpha|*}}_{\mathrm{DR}}(B/A)$ [NB $\Lambda\simeq\widetilde B/$\unsure{idéal|ident}] donc est un isom si $\widetilde B\xrightarrow{\sim}\Lambda$ i.e. $\widetilde B\longrightarrow B$ est injectif.

En effet, si $(x_i)_{i\in I}$ est une $p$-base, utilisant les notations du n$^\circ$ 1, on trouve \deleted{\unsure{un}}\inserted{un isom, de $B$-algèbres}
\diagram{\begin{tikzcd}[ampersand replacement=\&,column sep=1em,row sep=.75em]
 B\arrow[r]\arrow[d] \& \widetilde B\arrow[r]\arrow[d,dash] \& B\arrow[d]\\
 B_i\arrow[r,dash]\arrow[d] \& \widetilde B_i\arrow[r,dash]\arrow[d,dash] \& B_i\arrow[d]\\
 \Lambda\arrow[r,dash]\arrow[d,dash] \& \widetilde\Lambda\arrow[r,dash]\arrow[d,dash] \& \Lambda\\
 A\arrow[r,dash] \& \widetilde A \& {}
\end{tikzcd}}
\[
 \Omega^*_{B/A}\simeq\bigotimes_{i\in I}\Omega^*_{B_i/\Lambda}\otimes_{B_i}B
 \quad\text{d'où}\quad
 \Omega^\alpha_{\widetilde B/\widetilde A}\simeq\bigotimes_{i\in I}
 (\overset{\deleted{\illegible{1}}}{\Omega^*_{\widetilde B_i/\widetilde\Lambda}}\otimes_{\unsure{B_i|\widetilde B_i}}\widetilde B),
\]
\margin{isom de $B$-alg.}
et en vertu du calcul du n$^\circ$ 1 pour $H^\alpha_{\mathrm{DR}}(B/A)$ :
\[
 H^*_{\mathrm{DR}}(B/A)\simeq\bigotimes_{i\in I}H^\alpha_{\mathrm{DR}}(B_i/\Lambda)
\]
L'hom $\varphi$ envisagé s'identifie alors au produit tensoriel des homomorphismes analogues \deleted{\unsure{envisagés|correspondants}} de
\[
 \begin{aligned}
 \varphi_i:\Omega^*_{\widetilde B_i/\widetilde\Lambda}\otimes_{B_i}\widetilde B_i
 &\longrightarrow H^*_{\mathrm{DR}}(B_i/\Lambda)\\
 \text{base }\ \widetilde 1,\ d\widetilde x_i&\qquad
 \text{base }\ \dot{1},\ \overset{\cdot}{\overline{x_i^{p-1}dx_i}}
 \end{aligned}
\]
qui envoient base sur base, donc sont des isom. Cela prouve le théorème.

\uline{Corollaire} 1 \deleted{\unsure{Supposons de plus|Supplémentaire}} On a
\[
 Z^*(B/A)\simeq B^*(B/A)\oplus[\Omega^*_{\widetilde B/\widetilde A}\otimes_{\widetilde B}\Lambda]
\]

\uline{Définition}\quad L'isom inverse de $\varphi$ est noté
\[
 \boxed{\widetilde C:H^*_{\mathrm{DR}}(B/A)\xrightarrow{\sim}
 \Omega^*_{\widetilde B/\widetilde A}\otimes_{\widetilde B}\Lambda}
 \boxed{(=\Omega^\alpha_{B/\Lambda}\otimes_B\Lambda)}
\]
On dénote aussi souvent par $C$ l'hom.
\[
 \widetilde C:Z^*(B/A)\longrightarrow\Omega^*_{\widetilde B/\widetilde A}
\]
donnant le précédent par passage au quotient, [donc ce dernier est \uline{surjectif}, \uline{et son noyau est formé des différentielles exactes}].
% archive sheet number: 170

